%% =====================================================================
%%  T-28-01  Assume Formlessness
%%  -------------------------------------------------------------------
%%  One idea per file. Core is mandatory. Slots in canonical order.
%%  Max 700 words. No layout commands. No filler. New source material
%%  for this entry goes in sources/<BOOK>/T-28-01.tex -- never by
%%  rewriting this file (docs/RULES.md R0.1).
%%  =====================================================================
%% @kind: topic
%% @id: T-28-01
%% @title: Assume Formlessness
%% @subtitle: The final entry is the shape of every entry: never let them fix your form
%% @chapter: c28
%% @order: 10
%% @status: stable
%% @sources: B3
%% @updated: 2026-09-19
%% @allow-rewrite: B3 ingest 2026-09-19 -- committed scaffold replaced by the written entry (planned -> stable lifecycle)

\topic{T-28-01}{Assume Formlessness}{The final entry is the shape of every entry: never let them fix your form}

\begin{core}
A form is a target. Once your habits, plans and positions take a visible shape, opponents aim at the shape and society prices you by it --- the only way to be unattackable is to have no fixed form to attack. The final entry is a stance rather than a move: keep your strategy mobile, your beliefs unpredictable to observers, and your dependence on any single footing low --- let form be assumed only where form itself is the weapon.
\end{core}
\begin{mechanism}
Every fixed form offers its attacker three gifts: predictability, exposure and counter-formulation --- your shape tells the room where you will be, what you value, and how to be the opposite. The source's terminal image is the formless serpent: nothing in the animal is fixed --- scale by scale --- and precisely that mobility makes it untouchable; society attacks the grotesque (the openly deformed position, the declared obsession) and steers clear of the formless, because there is nothing to grab. The historical case is the officer with no uniform: operating amid enemies without insignia or fixed position, he was everywhere and nowhere --- his formlessness was his form. The stance scales down: refuse to become predictable in manoeuvre, keep your budget of exposure low, and reserve one face that no market has priced. The terminal discipline ties the book together: form is adopted as tactic, never lived as identity --- assume a form when it serves, shed it when it is priced, and be seen assuming it by no one whose instrument you would rather not become.
\end{mechanism}
\begin{conditions}
Strongest under observation --- negotiation, competition, politics, any room where your habits are being modelled. It weakens where trust requires legibility: teams, marriages, partnerships price your formlessness as unreliability, and the compliant spouse or the dependable colleague is the person whose form never moves. The stance is a professional posture, not a private one; practised at home it costs exactly what the intimacy chapters charge for concealment.
\end{conditions}
\begin{application}
  \item Audit your fixed forms: the route, the schedule, the stated position, the predictable counter. Each is a published target.
  \item Break your own patterns periodically for no reason. Unpredictability decays; assume no pattern stays unobserved past a season.
  \item Reserve one open channel --- skill, asset, allegiance --- that no employer, partner or audience can see or price.
  \item Adopt forms deliberately when they serve (the courtier's agreeableness, the bold entry), and shed them the season they are priced.
  \item With intimates, close the ledger: formlessness there is just distance. Choose the rooms where you will be knowable, and be knowable in them.
\end{application}
\begin{feedback}
  \item Working: opponents prepare for your last move and meet something else. Your shape never arrives on time for them.
  \item Working: your reserved channel lets you survive a loss that should have ended you.
  \item Failing: your own team cannot find you to promote you. Formlessness has eaten your legibility.
  \item Failing: observers name your pattern anyway --- \emph{you always} --- and price you. The stance needs a maintenance pass.
\end{feedback}
\begin{failure}
The stance's terminal cost is the terminal one: formlessness scales to the self, and the person with no fixed form eventually has nothing that persists --- no doctrine held past its pricing, no allegiance past its utility. The book's own last line is the confession: the serpent's advice is not for the decent --- power's final truth is a tools-not-people test, and the moment tools become the self, people become tools too. Practised as identity rather than posture, the stance delivers exactly that.
\end{failure}
\begin{countermeasures}
  \item Model others: does anyone in your life have no form you can map? Ask what that costs them, and what it costs you to trust them.
  \item Price the formless adversary by tracks, not shapes --- what moves, what gets funded, who benefits. Outcomes leave wakes.
  \item Demand legibility where you depend on someone; formlessness across the board is a negotiation position, not a personality.
  \item Keep your own core fixed on purpose --- one documented self, visible to the people who matter. Formlessness is armour for the surface, never for the core.
\end{countermeasures}

\sources{B3}
\seesources
\seealso{T-13-02, T-18-04, T-20-02}
